\newcommand{\tipoRequerimento}{Solicitação de reconhecimento e pagamento de valores complementares decorrentes de cobrança indevida}
\newcommand{\tese}{CORREÇÃO MONETÁRIA E JUROS DE PAGAMENTO}

\section{DOS FATOS E DOS DIREITOS}

\begin{enumerate}[resume]
    \item O Município, na condição de consumidor dos serviços de distribuição de energia elétrica, identificou a ocorrência de cobrança indevida referente à Unidade Consumidora \unidadeConsumidora{}.   

    \item Diante da constatação da irregularidade, o Município apresentou reclamação administrativa [RECLAMACAO], em [DATA RECLAMACAO], perante a concessionária requerendo as providências cabíveis para a restituição dos valores cobrados em desacordo com as normas regulamentares.

    \item Após a análise da reclamação, a concessionária reconheceu a cobrança indevida e o direito à restituição.

    \item Em decorrência desse reconhecimento, a concessionária realizou o primeiro pagamento em favor do Município em [DATA PRIMEIRO PAGAMENTO], no valor de R\$ [VALOR PRIMEIRO PAGAMENTO], conforme comprovante nº [COMPROVANTE PRIMEIRO PAGAMENTO].

    \item Ocorre que, após a análise do referido pagamento, o Município constatou que o valor efetivamente restituído não correspondia à integralidade do montante devido, em razão de interpretação equivocada ou erro no cálculo efetuado pela distribuidora.

    \item Diante da insuficiência do ressarcimento, o Município apresentou contestação demonstrando fundamentadamente o equívoco no montante devolvido.

    \item Após a análise do caso, Foi reconhecido o direito do Município à complementação dos valores indevidamente retidos.

    \item Em cumprimento à decisão, a concessionária realizou novo pagamento em favor do Município, de natureza complementar, no valor de R\$ [VALOR PAGAMENTO COMPLEMENTAR], efetuado em [DATA PAGAMENTO COMPLEMENTAR], conforme comprovante nº [COMPROVANTE PAGAMENTO COMPLEMENTAR].

    \item Contudo, o pagamento complementar solucionou apenas a parcela principal, não contemplando a devida atualização monetária e os juros de mora incidentes sobre o período em que o valor permaneceu indisponível para o ente público municipal.

    \item Verifica-se, portanto, que entre [DATA DISPONIBILIZACAO] e [DATA EFETIVO PAGAMENTO COMPLEMENTAR] transcorreu período de aproximadamente [PERIODO DECORRIDO], durante o qual o Município permaneceu privado da disponibilidade do montante reconhecidamente devido.

    \item Considerando que o pagamento do valor complementar ocorreu somente em [DATA PAGAMENTO], cumpre proceder à apuração dos efeitos financeiros decorrentes do período transcorrido entre [DATA TERMO INICIAL] e [DATA TERMO FINAL], especialmente no tocante à atualização monetária e aos juros de mora aplicáveis.

    \item Dessa forma, permanece controvérsia quanto à recomposição integral dos efeitos financeiros decorrentes do período em que o valor reconhecidamente devido permaneceu sem a devida disponibilização ao Município.

    \item A presente manifestação tem por finalidade submeter novamente a questão à apreciação da ANEEL, a fim de que seja analisada a existência do saldo financeiro decorrente do período de atraso e, sendo confirmado o direito, seja determinada à concessionária a respectiva complementação.
\end{enumerate}

\section{DOS PEDIDOS}

\begin{enumerate}[resume]
\item Diante do exposto, requer-se:
    \begin{enumerate}
        \item O recebimento e a análise da presente manifestação, vinculando-a, se cabível, ao Processo ANEEL nº [NUMERO PROCESSO ANEEL] e aos demais processos/protocolos relacionados ao caso;

        \item O reconhecimento da cronologia dos pagamentos realizados pela concessionária, especialmente:
        \begin{itemize}
            \item Primeiro pagamento, realizado em [DATA PRIMEIRO PAGAMENTO], no valor de R\$ [VALOR PRIMEIRO PAGAMENTO];
            \item Pagamento complementar, realizado em [DATA PAGAMENTO COMPLEMENTAR], no valor de R\$ [VALOR PAGAMENTO COMPLEMENTAR];
        \end{itemize}

        \item O reconhecimento de que o pagamento complementar realizado em [DATA PAGAMENTO COMPLEMENTAR] não contemplou a integralidade da atualização monetária e dos juros de mora devidos no período;

        \item A análise do período compreendido entre [DATA TERMO INICIAL] e [DATA TERMO FINAL] (intervalo entre [EXPLICACAO DATA INICIAL] e [EXPLICACAO DATA FINAL]), para fins de apuração e incidência da atualização monetária e dos juros de mora;

        \item A determinação do critério de atualização monetária e juros de mora aplicável ao caso, considerando a legislação e a regulamentação da ANEEL vigentes e aplicáveis ao período e à natureza da cobrança objeto da presente manifestação;

        \item A determinação à concessionária para realizar novo cálculo do valor devido, tomando como base o valor deferido inicial, o período de [DATA TERMO INICIAL] a [DATA TERMO FINAL] e os critérios definidos pela ANEEL;

        \item O reconhecimento da existência de saldo remanescente correspondente à atualização monetária e aos juros incidentes sobre o valor que permaneceu indisponível ao Município durante o período indicado;

        \item A determinação à concessionária para que efetue o pagamento do saldo remanescente, estimado pelo Município, ou no valor que vier a ser apurado pela ANEEL/concessionária segundo os critérios regulamentares aplicáveis;

        \item A apresentação, pela concessionária, de memória de cálculo detalhada, discriminando: valor principal considerado, termo inicial e final da atualização, índice de correção utilizado, percentual de juros aplicado, quantidade de dias considerada, valor da atualização monetária, valor dos juros e montante total devido;

        \item A apresentação dos comprovantes de pagamento e dos respectivos cálculos, permitindo ao Município verificar a integralidade do cumprimento da decisão da ANEEL e das obrigações decorrentes da cobrança indevida;

        \item Caso a ANEEL entenda não ser cabível o pagamento pleiteado, que seja apresentada manifestação fundamentada indicando expressamente a razão jurídica para o indeferimento e o critério utilizado para afastar a incidência da atualização monetária e/ou juros no período questionado.
    \end{enumerate}
\end{enumerate}